\textbf{Answer: 1) $T(y) = \rho g y$}

\underline{\textbf{Solution for 1}}

Let $T(y)$ be the tension as a function of height. Consider a small piece of the rope between $y$ and $y+dy$ ($0 \le y \le L$). The forces on this piece are $T(y+dy)$ upward, $T(y)$ downward, and the weight $\rho g\, dy$ downward. Since the rope is at rest, we have $T(y+dy) = T(y) + \rho g\, dy$. Expanding this to first order in $dy$ gives $T'(y) = \rho g$. The tension in the bottom of the rope is zero, so integrating from $y = 0$ up to a position $y$ gives

$$T(y) = \rho g y. \eqno(2.16)$$

As a double-check, at the top end we have $T(L) = \rho g L$, which is the weight of the entire rope, as it should be.

Alternatively, you can simply write down the answer, $T(y) = \rho g y$, by noting that the tension at a given point in the rope is what supports the weight of all the rope below it.